% SOURCE_STATEMENT: "Is potential-temperature inversion complementary to potential-vorticity inversion?"
% Preparation and final self-audit: CONTRIBUTING.md prompts C and D; source/audit artifacts retained in the local physics_statements_2026_10_11 batch.
% !TeX program = lualatex
\documentclass[11pt]{article}
\usepackage[margin=25mm]{geometry}
\usepackage{fontspec}
\setmainfont{Latin Modern Roman}
\usepackage{amsmath,amssymb,mathtools}
\usepackage{unicode-math}
\setmathfont{Latin Modern Math}
\usepackage{xurl,hyperref}
\hypersetup{colorlinks=true,urlcolor=blue}
\setcounter{secnumdepth}{0}
\setlength{\parindent}{0pt}
\setlength{\parskip}{5pt}
\setlength{\emergencystretch}{3em}
\newcommand{\sref}[2]{\hyperlink{source:#1:#2}{[S#2]}}
\newcommand{\cataloguescope}[1]{\paragraph{Recorded target.}#1}
\newcommand{\R}{\mathbb R}
\newcommand{\C}{\mathbb C}
\newcommand{\N}{\mathbb N}
\newcommand{\Z}{\mathbb Z}
\newcommand{\Q}{\mathbb Q}
\newcommand{\F}{\mathbb F}
\newcommand{\eps}{\varepsilon}
\begin{document}
% UnsolvedMath ID 30001530; source catalogue code OWR-4415-003
\setcounter{section}{30001529}
\section{Potential Temperature and Potential Vorticity Inversion}\label{30001530}
{\small UnsolvedMath ID 30001530 · Mathematical Physics}
\subsection{Definitions and mathematical statement}

\paragraph{Shared notation.}$\mathbb N=\{1,2,\ldots\}$; the standard number fields have their usual meanings. All additional conventions are stated locally.


\paragraph{Definitions and conventions.} Potential temperature is the temperature a parcel would have after adiabatic transfer to a reference pressure; in dry ideal-gas notation, $\theta=T(p_{\rm ref}/p)^{R/c_p}$. Potential vorticity combines absolute vorticity and a stratification gradient; Ertel's notation is $q=\rho^{-1}(\boldsymbol\omega+2\boldsymbol\Omega)\cdot\nabla\theta$. Under a chosen balanced-flow relation, potential-vorticity inversion recovers other fields from $q$ and boundary data. Potential-temperature inversion starts from $\theta$ and hydrostatic/balance conditions to recover pressure, winds, and vorticity.

\paragraph{Source scope and qualifications.} Complementarity concerns information and diagnostic interpretation under consistent balance and boundary assumptions, not unconditional unique recovery from either scalar field. The report's simplest PT inversion uses hydrostatic integration with a boundary condition; the boundary information cannot be discarded.

\paragraph{Catalogue mathematical target.}
\paragraph{Catalogue wording retained for traceability.}
Is potential-temperature inversion complementary to potential-vorticity inversion?

\paragraph{Saved literature context.} The supplied catalogue's assessment dated 2026-08-21 records: The report presents potential-temperature inversion as a less-developed alternative to potential-vorticity inversion and leaves the complementarity question open; no general theorem settling it was identified. The assessment is retained as dated provenance; it is not a new comprehensive literature or open-status certification.
\subsection{Short English statement}
Determine how inversion from potential temperature complements the more usual inversion from potential vorticity.
\subsection{Sources}
\begin{itemize}
\item[\hypertarget{source:30001530:1}{S1}] ULAM.AI and the UnsolvedMath Contributors, UnsolvedMath, record 30001530 (OWR-4415-003). Catalogue curation under CC BY 4.0; definitions and English exposition added locally.
\url{https://huggingface.co/datasets/ulamai/UnsolvedMath}.
\item[\hypertarget{source:30001530:2}{S2}] Mathematical Theory and Modelling in Atmosphere-Ocean-Science (2010), Oberwolfach Reports, DOI 10.4171/owr/2010/34, pp. 2027--2099; official report source and mathematical context.
\url{https://doi.org/10.4171/owr/2010/34}.
\item[\hypertarget{source:30001530:3}{S3}] Official EMS Press PDF of the same Oberwolfach report; PDF pages 18--19, printed pages 2044--2045.
\url{https://ems.press/content/serial-article-files/46292}.
\end{itemize}
\label{end:30001530}
\end{document}
